\chapter{Conclusiones}
\label{ch:conclusiones}

El desarrollo del Ejercicio Práctico N° 3 ha permitido validar empíricamente las ventajas de migrar
un desarrollo \textit{Bare Metal} puro hacia el estándar de abstracción de hardware CMSIS en microcontroladores
de arquitectura ARM Cortex-M.

Se logró portar la totalidad del código fuente —incluyendo la inicialización del reloj, la configuración
de los puertos GPIO, la parametrización de interrupciones externas (EXTI/NVIC) y la gestión del
convertidor analógico-digital (ADC1)— sustituyendo el uso riesgoso de punteros directos a memoria
por el empleo sistemático de estructuras \lstinline{struct} tipadas e identificadores mnemotécnicos. 
Esta transición no solo previene errores de direccionamiento originados por errores en el cálculo
de los desplazamientos (\textit{offsets}), sino que mejora significativamente la legibilidad, escalabilidad
y mantenimiento del firmware.

Adicionalmente, la implementación del temporizador \textit{SysTick} permitió erradicar los retardos
empíricos por bucles bloqueantes de software. La base de tiempo precisa y determinista de \qty{1}{\ms}
alcanzada demuestra la potencia de integrar periféricos internos del núcleo para tareas de temporización.

Como grupo, las principales dificultades encontradas radicaron en la configuración inicial del \textit{SysTick}
—específicamente el entendimiento de la relación entre la frecuencia del sistema y el registro de recarga— y, de manera
más general, en la transición desde el uso de desplazamientos de bits con valores numéricos manuales hacia el empleo de
las macros semánticas provistas por el estándar. Este último proceso resultó laborioso, ya que requirió una búsqueda
exhaustiva en la documentación de la biblioteca para identificar la constante exacta correspondiente a cada campo de
registro, reemplazando la lógica de "números mágicos" por identificadores estandarizados. 
